% Transcription — fonds Alexandre Grothendieck, shelfmark 131, batch 1 (pages 1-20).
% Established from the facsimile archives/batches/131/batch-01.pdf (a mirror of
% https://grothendieck.umontpellier.fr/131.pdf), per the skill
% transcribe-grothendieck. The batch file carries no cover sheet: PDF page N
% is page N of the folder, and the archivists' pencil numbers bottom-left
% agree wherever legible (« 1 », « 3 », « 4 », « 6 », « 8 », « 10 », « 11 »,
% « 13 », « 15 », « 16 », « 17 »).
%
% Pass: Opus 5.5 (claude-opus-5-5), 2026-09-24 — first pass, unchecked against the pages by a human.
%
% Dating. The inventory dates the folder « [vers 1971] » (copied verbatim
% into \dating). The folder sits in the inventory group « Descentes »
% (126-128, which also takes 129-133), dated « [avant 1970] »: the group's
% date and the folder's own do not agree. Recorded as an observation only.
%
% Fourteen of the twenty pages are transcribed: 1, 2, 3, 4, 5, 6, 8, 11,
% 13, 14, 15, 16, 17, 19 (two of them, 5 and 17, a single line). Absent:
% pages 10 and 18, blank; pages 7, 9, 12 and 20, leaves of a typed text by
% someone else (headed « n° 392 », pages « 4-27 », « 4-26 », « 4-24 », and
% « n° 371 », page « -34- »: classes of divisors attached to modules of
% finite type, relative invariants of lattices, graded algebras), with no
% mark in his hand and no bearing on the mathematics of the folder —
% skipped under the project's rule for other people's text. Pages 6, 8 and
% 19 are written on the backs of such typed leaves (their text shows
% through); only his hand is transcribed. He did not paginate these leaves;
% page 1 carries a circled « A » top right.
%
% What the batch is. A set of conditions on a functor F : (Sch/S)° → Ens,
% S locally noetherian, for F to be representable by a scheme locally
% quasi-finite and separated over S, and the proof by induction on dim S.
%   1      The eight conditions (fpqc sheaf; commutes with filtered direct
%          limits; commutes with adic inverse limits of artinian rings;
%          left exactness on finite algebras over an artinian ring;
%          pro-representability by finite algebras at points with values
%          in a field; separation by the valuative criterion; « critère de
%          modularité » local (7°) and generic (8°)).
%   2      Pencil: what to replace for unramified functors.
%   3-4    The proof: stability under base change; (1) dim S = 0; (2) S
%          local complete of dimension n > 0, by induction on the punctured
%          spectrum, using « prop. 2 des notes » (earlier notes, not in this
%          batch); (3) S local, by descent from the completion; (4) S
%          arbitrary, by a lemma reducing to the local rings (« prop. 3 des
%          notes préc. »).
%   5      One pencil line (constructibility, « loc. cit. »).
%   6, 8   Reformulations of 7° and 8° with nilpotent thickenings, and the
%          axioms (A), (B), (C) (open immersion at a point; openness of the
%          étale locus, C_1 stable under generisation, C_2); a first proof
%          of (A) on page 19.
%   11     A second list of the conditions, numbered 1° to 9° (a clean copy
%          in progress, with a column of dense vertical notes in the left
%          margin, largely illegible).
%   13-16  The « Théorème » stated in full (conditions 1° to 8°), N.B. on
%          which conditions are needed; Corollaire 1 (étale and separated),
%          Corollaire 2 (unramified and separated), Remarque 3 on 8°;
%          condition (RG); Remarque 4; « Immersion ouverte » (page 14);
%          examples of unramified functors and « Applications en vue »
%          (projectivity of abelian schemes, existence of Pic).
%   17     « Façon d'exprimer », alone on the leaf.
%   19     Proof of (A) and (B): X ×_{F_{S'}} X is a closed subscheme of
%          X ×_S X formally equal to the diagonal, so X → F is a
%          monomorphism, étale, hence an open immersion.
%
% Notation fixed for the batch. \xi for the curly letter, shaped like a
% reversed 3, that names points and morphisms (ζ : Z → F_{S'}, ξ ∈ F(k));
% the page does not distinguish a ζ from a ξ and one letter is kept. Z for
% the doubled, bold-looking Z of page 1 (a finite S'-scheme). S'_0 for his
% « S'_0 ». « tt », « pt », « loc. », « qu. », « noeth. », « fid. »,
% « imm. », « V.A. », « néc. », « hyp. » are kept. The conditions are
% numbered 1°, 2°, … as he writes them; the numbering differs between
% pages 1, 11 and 13 and is not harmonised. Square brackets in the text are
% his.
%
% The dating is the inventory's own for the folder, copied verbatim; the
% title in \foldertitle is the archivists' (not bracketed in the inventory).
\documentclass[11pt,a4paper]{article}
% Le préambule commun à toutes les transcriptions du fonds Grothendieck.
%
% Il tient en une idée : **l'incertitude se compose, elle ne se résout pas.**
% Une transcription qui lisse un mot douteux a détruit la seule chose pour
% laquelle on la faisait — un lecteur doit pouvoir savoir, à chaque endroit,
% ce qui était sur le feuillet et ce qui a été ajouté. Les six macros qui
% suivent sont donc l'appareil critique, et non de la décoration.
%
% Les mêmes commandes sont reconnues par `scripts/render.mjs`, qui produit la
% vue de lecture du site. Une macro ajoutée ici doit l'être aussi là-bas,
% faute de quoi le rendu échouera bruyamment — ce qui est voulu : un rendu
% silencieusement faux est pire qu'un rendu absent.

\ProvidesPackage{grothendieck}[2026/08/08 Transcriptions du fonds Grothendieck]

\usepackage{amsmath,amssymb,amsthm}
\usepackage{mathrsfs}
\usepackage{stmaryrd}
% Les diagrammes commutatifs sont partout dans ce fonds : c'est la langue
% même de la géométrie algébrique de Grothendieck, et une transcription qui
% les rendrait en prose ne transcrirait rien.
\usepackage{tikz-cd}
% Français par défaut — c'est la langue des feuillets — mais l'anglais est
% chargé aussi : la traduction et le résumé sont des documents anglais, et la
% typographie française (espaces avant les deux-points) y serait une faute.
% Ces documents basculent par \selectlanguage{english} en tête de corps.
\usepackage[english,french]{babel}
\usepackage{fontspec}
\usepackage{geometry}
\geometry{a4paper,margin=25mm,marginparwidth=28mm}
\usepackage{marginnote}
\usepackage{xcolor}
% ulem, not soul. `soul`'s \st and \ul are built on a character-by-character
% scan that XeLaTeX's font machinery breaks: the document fails with an
% undefined control sequence rather than degrading. ulem does the same two jobs
% and is Unicode-engine safe. `normalem` stops it hijacking \emph, which the
% transcriptions use for Grothendieck's own emphasis.
\usepackage[normalem]{ulem}
\usepackage{hyperref}
\hypersetup{colorlinks=true,linkcolor=black,urlcolor=black,pdfborder={0 0 0}}

% --- Métadonnées du lot -----------------------------------------------------
% Reprises telles quelles par le script de rendu, qui les affiche en tête de
% la vue de lecture. Elles fixent ce que le fichier prétend couvrir : sans
% elles, une transcription flotte sans point d'attache dans un fonds de seize
% mille feuillets.

\newcommand{\folder}[1]{\gdef\grothFolder{#1}}
\newcommand{\batch}[1]{\gdef\grothBatch{#1}}
\newcommand{\leaves}[2]{\gdef\grothFirst{#1}\gdef\grothLast{#2}}
\newcommand{\foldertitle}[1]{\gdef\grothFoldertitle{#1}}
\gdef\grothFolder{?}\gdef\grothBatch{?}\gdef\grothFirst{?}\gdef\grothLast{?}\gdef\grothFoldertitle{}

% --- Le résumé --------------------------------------------------------------
%
% Il ouvre la lecture modernisée, sous le titre « Résumé », et il est composé
% en petit. Ce n'est pas un ornement : c'est le seul passage du document qui
% ne soit pas des mathématiques, et un lecteur doit pouvoir le reconnaître
% avant de l'avoir lu — puis le sauter s'il connaît déjà le sujet, ou s'y
% arrêter s'il ne le connaît pas. Le filet qui le suit marque la reprise du
% corps.

\newenvironment{resume}
  {\small\setlength{\parskip}{0.45em}}
  {\par\medskip\hrule\medskip}

% --- Mots-clés ---------------------------------------------------------------
%
% En anglais, en clôture du résumé de la lecture modernisée : le vocabulaire
% moderne sous lequel chercher ce que les feuillets construisent. Le manifeste
% du site (`npm run manifest`) les extrait pour étiqueter la cote entière —
% cette ligne est la source unique des tags, il n'y a pas de fichier à côté.
\newcommand{\keywords}[1]{\par\smallskip\noindent
  {\footnotesize\textsc{Keywords} --- \textit{#1}}\par}

% --- L'appareil critique ----------------------------------------------------

% Le numéro de feuillet, en marge. C'est l'ancre qui permet de régler un
% désaccord sur une formule en regardant *un* feuillet plutôt qu'une chemise
% de six cents. Le site s'en sert aussi pour tourner les pages du fac-similé
% quand on fait défiler la transcription.
\newcommand{\leaf}[1]{%
  \marginnote{\footnotesize\color{gray!70}\textsf{#1}}%
  \label{leaf:#1}\ignorespaces}

% Illisible : on ne devine pas. Un mot inventé qui se lit comme les autres est
% le pire résultat possible.
\newcommand{\ill}{\textcolor{red!60!black}{[\,\ldots\,]}}

% Lecture douteuse : le mot est donné, et signalé comme incertain.
\newcommand{\uncertain}[1]{\uline{#1}}

% Ajout de l'éditeur — une lettre manquante, un mot sous-entendu.
\newcommand{\add}[1]{\textcolor{blue!50!black}{[#1]}}

% Biffé par Grothendieck lui-même. Ce qu'il a rayé fait partie du document :
% une rature dit souvent où la pensée a changé de direction.
\newcommand{\struck}[1]{\sout{#1}}

% Note du transcripteur, sur l'état matériel du feuillet.
\newcommand{\note}[1]{\par\smallskip{\small\color{gray!60!black}\textit{[#1]}}\par\smallskip}

% Note marginale de Grothendieck, distincte du corps du texte.
\newcommand{\marginal}[1]{\par\smallskip\noindent\hspace{1em}%
  {\small\color{gray!60!black}#1}\par\smallskip}

% --- Titre ------------------------------------------------------------------

\AtBeginDocument{%
  \begin{center}
    {\large\bfseries Fonds Alexandre Grothendieck --- cote n\textsuperscript{o}~\grothFolder}\\[2pt]
    {\itshape\grothFoldertitle}\\[4pt]
    {\small feuillets \grothFirst--\grothLast\ (lot \grothBatch)}\\[2pt]
    {\footnotesize\color{gray!60!black}Transcription établie d'après le fac-similé numérisé
     par l'Université de Montpellier.\\
     Les lectures incertaines sont soulignées, les illisibles notées [\,\ldots\,],
     les ajouts entre crochets.}
  \end{center}
  \vspace{1em}}

\endinput


\folder{131}
\batch{1}
\pages{1}{20}
\dating{[vers 1971]}
\watermark{Édition de démonstration}
\foldertitle{Foncteurs représentables par schéma quasi fini : notes manuscrites (s.d.).}

\begin{document}

\page{1}
\note{en haut à droite, un « A » entouré ; les conditions 1° à 4° sont
réunies par une accolade à l'encre, et reliées au crayon, par des traits
qui convergent dans la marge, à la mention « \uncertain{conditions}
\uncertain{d'exactitude} » et aux signes $\varprojlim$, $\varinjlim$}

1° $F$ faisceau pour la top. fpqc
$\left\{\begin{array}{l}
\text{a) de nature locale Zariski} \\
\text{b) compatible avec descente fid. plate qu. cpte affine}
\end{array}\right.$
\note{« de nature locale Zariski » est écrit à l'encre brune, comme
« affine » en fin de ligne}

2° $F$ commute aux $\varinjlim$ d'anneaux (propriété \emph{relative})

3° $F$ commute aux $\varprojlim$ adiques d'anneaux \emph{artiniens}
\note{« artiniens » est souligné trois fois}

4° Pour tout anneau artinien $A$ sur $S$, $F_A$ est exact i.e. sur la
catégorie des algèbres \struck{\ill{}} \struck{artiniennes} finies sur $A$
[\uncertain{on a ainsi} une suite exacte
\[
  A \to A' \rightrightarrows A' \otimes_A A' = A''
\]
$A$ artinien, local \struck{\ill{} artiniens}, avec $A'$ \struck{\ill{}}
finie sur $A$, [ext. résiduelles triviales,
$\mathrm{long}_A\, A'/A = 1$] donne une suite exacte
\[
  F(A) \to F(A') \rightrightarrows F(A'')
\]
].

\marginal{finitude ou discrétion}
\marginal{peut aussi se formuler pour $X/k$}
\note{deux notes au crayon dans la marge de gauche, en face de 5°, la
première soulignée}

5° \struck{\ill{}} Pour tout anneau \struck{\ill{}} \struck{noeth.}
\uncertain{corps} \ill{} $k$ sur $S$, et $\forall\, \xi \in F(k)$,
l'anneau $B$ sur $k$ qui pro-représente $F_k$ en $\xi$ est \emph{fini} sur
$k$. (propriété \emph{relative})
\note{la correction au-dessus de la ligne, où se lit « corps », remplace
un premier « anneau local » biffé ; le $k$ est surchargé}

\marginal{$F(k) \to F(X)$ \ill{} \ill{} (\ill{} mod $k$) se réduit comme
\ill{} (cas : $X$ \uncertain{intègre} \ill{})}
\note{note au crayon, oblique, écrite en travers de la ligne de 5°, entre
le numéro et le texte ; la flèche de $F(k)$ vers $F(X)$ est double ou
barrée, lecture douteuse}

\marginal{séparation}

6° $F$ est \emph{séparé} sur $S$ (propriété \emph{relative})
(par critère valuatif)

\marginal{critère de modularité : a) locale}

7° Pour tout $S'$ local noeth. complet [sur $S$], tout $Z$ [local] fini
sur $S'$, pt fermé $t$, \struck{\ill{}} tout $\xi : Z \to F_{S'}$
[i.e. $\xi : Z \to F$] qui est modulaire (pour $F_{S'}$) au pt fermé $t$ de
$Z$, $\xi$ est modulaire [ou tout au moins, étale] en tout point.

\begin{tikzcd}
  F \arrow[d] & F_{S'} \arrow[l] \arrow[d] & \\
  S & S' \arrow[l, no head] & Z \arrow[l] \arrow[ul, "\xi"']
\end{tikzcd}

\note{le diagramme est dessiné à gauche du texte de 7° ; le trait entre
$S$ et $S'$ n'a pas de pointe ; les deux traits verticaux descendent de
$F$ vers $S$ et de $F_{S'}$ vers $S'$}

\marginal{On peut \uncertain{prendre} $S'$ irréductible de dim 1 (et
$F_{S'\ill{}}$ représentable, \ill{} \struck{\ill{}} \ill{} fermé de $S'$)
$\to$ 8°}
\note{dans la marge de gauche, entre 7° et 8°, à l'encre, en partie
coupé par une accolade au crayon ; une flèche renvoie à 8°}

\marginal{b) générique}

8° Pour tout $S'$ noeth. irréductible sur $S$, et $Z$
[\uncertain{irréductible}] \struck{$S'$} fini sur $S'$,
$\xi : Z \to F_{S'}$, tel que $\xi$ soit modulaire pour $F_{S'}$ au pt
générique $z$ de $Z$, il existe un ouvert $U \ni z$ tel que $\xi$ soit
modulaire (ou tt au moins étale) en tt point de $U$.
\note{« irréductible » est écrit au crayon au-dessus de $Z$, avec une
accolade ; « ou tt au moins étale » est interlinéaire}

\page{2}
\note{page au crayon, sans titre}

Dans le cas des critères pour les foncteurs non-ramifiés remplacer 5° par
la non ramification (\uncertain{en outre}) et \ill{} 4° et 5°,
\ill{} une condition \ill{} aussi simple que \uncertain{pour la}
[construction des foncteurs pro-représentables \emph{non ramifiés}]

\page{3}
On notera que les conditions 1° à 8° sont stables par toute extension de
la base. Sauf 2°, 5°, 6°, ce sont des conditions \emph{locales} sur $F$.

Pour prouver la représentabilité de $F$, on commence par le faire pour
$\dim S < +\infty$, \struck{\ill{}}

(1) $\dim S = 0$. On est ramené au cas $S$ local artinien, on construit un
\struck{préschéma} loc. qu. fini sur $S$ grâce à : 4°, 5°, on voit qu'il
représente $F$ grâce à : 1°, 2°, 3° [cf. notes précédentes, cor 2 à
prop. 1] OK. NB On n'a eu besoin que de 1° à 5°.
\note{les numéros de conditions de cette page sont entourés au crayon}

\marginal{N.B. Si $\dim S = 1$ ($S$ local complet), l'hypothèse 8° est
inutile.}

(2) $\dim S = n > 0$, $S$ local \emph{complet}. Soit $s$ son pt fermé, et
$U = S - s$, alors par l'hypothèse de récurrence [où 1° à 5° suffisent si
$S$ de dim 1] on sait que $F_U$ est représentable [par un $X_U$]. De plus,
grâce à : 2°, 5°, on construit une famille (utilisant $S$ complet) de
schémas \emph{finis} $X_i$ sur $S$, correspondant aux pts de $F_0$, et
\struck{\ill{}} grâce à : 3° des homs $u_i : X_i \to F$. En vertu de la
prop. 2 des notes, on est ramené à vérifier que
$u_i|U : X_{iU} \to X_U$ sont des \struck{immersions} immersions
\emph{ouvertes}, à images deux à deux disjointes. Ce dernier point provient
de l'hypothèse
\note{« $X_U$ » est écrit au-dessus du premier « immersions », biffé}

\page{4}
\marginal{\uncertain{Dégager} un \uncertain{lemme} \ill{} effet}
\note{note à l'encre brune dans la marge de gauche, à la hauteur des
lignes 3 et 4}

6°. Utilisant la condition 7°, on est ramené à montrer que
$X_i \xrightarrow{u_i} F$ est modulaire en le centre $x_i$ de $X_i$. Comme
la restriction $u_{i0} : X_{i0} \to F_0$ est modulaire (donc
\struck{\ill{}} le $k(x_i)$ est un corps de définition pour
$u_i : k(x_i) \to F$ sur $S$] il reste à vérifier la condition modulaire.
Nous \ill{} \ill{} d'abord démontrer le théorème dans le cas où la
condition 7° est \uncertain{renforcée}, en y exigeant seulement $\xi$
pro-\emph{modulaire} sur $S$ en $z$. De plus, $X$ est loc. qu. fini sur
$S$ par construction, et séparé sur $S$ par 6°.
\note{le « 6° » entouré en tête de page ne désigne pas la condition 6°,
mais la suite de l'hypothèse sur laquelle s'achève la page 3 ; le « 7° »
de la première ligne est entouré au crayon}

\marginal{Pour l'instant on n'utilise \uncertain{que} \struck{la forme
plus} \uncertain{une version affaiblie} \uncertain{mentionnée} de 7°}
\note{note oblique dans la marge de gauche, à la hauteur de la fin de
(2)}

(3) $\dim S = n > 0$, $S$ local. On sait par 2) que $F_{\hat{S}}$ est
représentable par un $X_{\hat{S}}$ loc. quasi-fini et séparé sur
$\hat{S}$. \struck{Il} Il en résulte que la donnée de descente
\uncertain{canonique} sur $F_{\hat{S}}$ est effective, \struck{Utilisant}
d'où un $X$ sur $S$, loc. quasi-fini sur $S$ et séparé sur $S$. Utilisant
1°, b) on trouve qu'il représente $F$.

(4) $\dim S = n$, $S$ quelconque. \struck{Comme en} On utilise le
[$F$ comme dessus]

\emph{Lemme}. Supposons que pour tout $a \in S$, $F_{S^a}$ sur
$S^a = \mathrm{Spec}\, \mathcal{O}_{S,a}$ soit représentable. Alors $F$
est représentable.

On utilise la \uncertain{prop.} 3 des notes préc.

\page{5}
\note{une seule ligne, au crayon, en haut de la page}

Compte tenu des conditions 1° et 2°, on est ramené à vérifier la condition
de constructibilité de loc. cit.

\page{6}
$C_1$ peut aussi se mettre sous la forme

7°) Pour tt $S' \to S$ [avec $S'$ \struck{local} noeth. [local]
\uncertain{complet} \uncertain{irréductible} de dim 1] et $S'_0$
sous-\struck{préschéma} schéma défini par un idéal nilpotent [de carré
nul], \struck{toute} ouvert $U \subset S'$ [$U = S' - $ pt fermé], et
\struck{section $\xi$ de $F_T$} $S'$-morphismes
$\xi_0 : S'_0 \to F_{S'}$, \struck{tel que et} \struck{$\xi : S'$} si
$\xi_0 | U_0$ se prolonge en $\xi_U : U \to F_{S'}$, alors $\xi_0$ se
prolonge en $\xi : S' \to F_{S'}$.

\marginal{d'adhérence schématique $S'$}
\note{note entourée dans la marge de gauche, reliée par un trait à
« $U \subset S'$ »}

[NB. Cette condition est \emph{vide} si $F \to S$ est étale]

\marginal{Ceci se met sous la forme :}
\note{au crayon, entouré, dans la marge de gauche}

8°) Soit $S'$ sur $S$ noeth. irréductible, [de pt gén. $x$] $S'_0$ défini
par un idéal \struck{nilpotent}, $\xi : S'_0 \to F_{S'}$, supposons que
$\xi$ soit modulaire [i.e. étale] en $x$, alors il est modulaire en un
voisinage de $x$.

Revenons au cas général, [(pas néc. non ram.)] je dis que (A), et (B)
peuvent se remplacer par

(D) (D${}_1$) Si $S'$ noeth. sur $S$, \struck{$X$} \struck{\ill{}}
$X$ \uncertain{fini} sur $S'$, et $\xi : X \to F_{S'}$ un $S'$-morphisme,
\ill{} $U$ \ill{} $\xi$ est \uncertain{inj.} étale est \emph{ouvert} ;
\struck{Et} $F \xrightarrow{\mathrm{diag}} F \times_S F$ (puis \ill{})
\ill{} il en est de même \ill{} \ill{} $X \times_{F_{S'}} X$ sur $S'$
\ill{} de sorte que cette condition \ill{} 7° et 8°) s'expriment
\ill{}.
\note{le bas de la page est un palimpseste : deux lignes au crayon
(« conditions \ill{} de base \ill{} 7° et 8° \ill{} ») biffées, un
cadre, une flèche vers $X \times_S X$, et plusieurs traits obliques qui
biffent ou relient ; seuls les mots rendus ci-dessus se lisent}

\page{8}
\emph{Axiomes} 1° à 6°, et

(A) Si $S'$ local noeth. complet [de dim 1 si on veut], la
\struck{pro}modularité de $X \xrightarrow{\xi} F_{S'}$ en $z$ sur le pt
fermé de $S'$ ($X$ fini sur $S'$) implique \uncertain{que} $\xi$
\struck{modulaire} imm. ouverte (\uncertain{Cela} \uncertain{si}
$F|(S' - \{s'\})$ repr.)

(B) Si \struck{aux des} $S'$ noeth. sur $S$, \struck{$X$} \ill{} sur $S$,
$\xi : X \to F_{S'}$ et $x \in X$ sont tels que $\xi$ \struck{modulaire}
\ill{} en $x$, alors $\exists\, U \ni x$ voisinage ouvert, t.q.
$\xi | U : U \to F_{S'}$ soit une imm. ouverte [On veut dire
\uncertain{simplement} que
$X \times_{\mathrm{Sp}\, \mathcal{O}_{S',s'}} \to F_{\mathrm{Sp}\, \mathcal{O}_{S',s'}}$
est une imm. ouverte]

\emph{Cas particuliers}

\struck{(i) $F \to S$ un \emph{monomorphisme}. Alors (A) et (B)}
contiennent les

(C) Pour tt $S'$ noethérien /$S$, \struck{\ill{}} et
$\xi : X \to F_{S'}$, [$X \to S'$ un \uncertain{mor.}], l'ens des
$x \in X$ en lesquels $\xi$ est [inj.] étale est ouvert.
\note{un long trait oblique traverse ces lignes, de « (i) » à
« l'ens des », et la condition (C) est encadrée ; on ne sait si le trait
biffe le passage}

\emph{Lemme} $\to$ On peut supposer en A resp. B que $X \to S'$ \uncertain{est}
une imm. \emph{fermée}. Les conditions (A) et (B), si $F \to S$ non
ramifié, se réduisent donc

\marginal{définie si on veut par un Idéal de carré nul}
\note{note de la marge de gauche, reliée à « fermée » ; la même mention
est répétée, entourée, à droite de la ligne suivante}

(C) Si $S'$ noeth. sur $S$, $S'_0 \to S'$ \emph{immersion fermée}
\struck{\ill{}} sur $S'$, $\xi : X \to F_{S'}$, l'ens. $U$ des pts en
lesquels $\xi$ est \emph{inj. étale} est \emph{ouvert} i.e.

(C${}_1$) stable par générisation [est aussi équivalent : A]

(C${}_2$) si $x \in U$, alors \struck{\ill{}} $U \cap \bar{x}$ est un
voisinage de $x$ dans $\bar{x}$.

\emph{N.B.} La condition (C) est automatique si $F \to S$ est
\emph{étale}.

\page{11}
\note{feuillet plus court ; la marge de gauche porte deux colonnes de
notes écrites verticalement, serrées et surchargées, que la présente
lecture ne restitue pas (on y distingue « Alors », $S'$, $\bar{X}$,
$S$, « N.B. », « 8° et 9° »), et, en bas à gauche, une note oblique sur
4° (« peut se remplacer par \ill{} : $A \to A'$ \ill{} $A'$ fini sur $A$
\ill{} [local] \ill{} $\mathrm{long}\, A'/A = 1$] alors
$F(A) \rightrightarrows F(A')$ »)}

(9°) Si $T$ noeth. irréd. sur $S$, $\exists\, U \subset T$ ouvert non vide
tel que $F_U$ soit représentable
\note{le numéro entouré, « 9° », est écrit par-dessus un autre, peut-être
« 10° »}

\marginal{peut se formuler de façon un \ill{} plus faible, avec
\uncertain{impossibilité} \ill{} de prolongement}
\note{note à l'encre brune, entourée, reliée par une flèche au numéro
« 9° »}

(1°) Foncteur $F$ de nature locale

(2°) $F$ compatible avec descente fid. plate qu.-compacte

(3°) $F$ commute aux $\varinjlim$ d'anneaux

(4°) $F$ commute aux $\varprojlim$ d'anneaux \uncertain{artiniens}
(i.e. les conditions \uncertain{habituelles} \uncertain{cela suffit} sur les
artiniens)

(7°) $F$ « non ramifié » \struck{ét} [$\ill{} = : \ill{}\ k[t]/(t^2)$]
\note{le « 7° » entouré est relié par un trait à la marge de gauche}

\struck{(6°) Tout pt de $F$ à valeurs dans une \struck{ext.} extension $K$,
\struck{est} sur $K$, provient d'une sous-extension finie}
\note{le numéro est barré en croix et les deux lignes sont biffées de
traits obliques au crayon}

(5°) $F$ pro-représentable universellement, \struck{\ill{}}
\note{à la suite, un cadre entièrement couvert de hachures}

(\uncertain{6°}) $F$ « séparé » (critère valuatif).

(8°) Pour tout $A$ anneau local noeth. de dim 1, [\uncertain{complet}]
irréductible sur $S$, et tt idéal $I$ de $A$ nilpotent, [i.e.
\ill{} $\mathfrak{p}$ idéal premier] \struck{et} $\xi \in F(A/I)$, si
$\xi_{\mathfrak{p}} \in \mathrm{Im}\, F(A_{\mathfrak{p}})$
\struck{alors $\xi \in \mathrm{Im}\, F(A/I)$} \struck{\ill{}}
\note{les deux dernières lignes, où se
lisent encore $\xi \in F(A/I)$, $\mathfrak{p} \in I$ et
$\xi_{\mathfrak{p}} \in F(A_{\mathfrak{p}}/I_{\mathfrak{p}})$, sont biffées
et reprises par une accolade ; une note entourée, sous le second
« 8° », est illisible}

\page{13}
\emph{Théorème}. Soit $S$ préschéma loc. noeth.,
$F : (\mathrm{Sch})_{/S}^{\circ} \to \mathrm{Ens}$. Pour que $F$ soit
représentable par un $X$ localement \emph{quasi-fini} sur $S$ et
\emph{séparé} sur $S$, il faut et il suffit que $F$ satisfasse les
conditions suivantes :

1°) $F$ de nature locale

3°) $F$ « commute aux $\varinjlim$ d'anneaux »

2°) $F$ compatible avec descente fid. plate et qu. compacte [il suffit
pour $B \to \hat{B}$, où $B$ est un anneau local noeth. sur $S$]
\struck{\ill{} de type fini et \ill{} cat. résiduelle \ill{} finie\ill{}]}
\note{« noeth. » est écrit au-dessus de « local » ; l'ordre 1°, 3°, 2°
est celui de la page}

4°) $F$ commute aux $\varprojlim$ d'anneaux artiniens, pour
$\hat{B} = \varprojlim B_n$, avec $B$ comme dans 2°

5°) $F$ pro-représentable au dessus de \struck{tt}
[tout anneau artinien $A$ sur $S$] \struck{\ill{}}, les anneaux $B$ qui
\struck{représentent} pro-représentent $F$ sur $A$ sont finis sur
\struck{$A$} $A$ ]
\note{la mention « anneau artinien $A$ sur $S$ » est entourée et reportée
d'une ligne plus bas par un trait}

6°) $F$ « séparé » (par critère valuatif …)

7°) Pour tout $A$, \struck{\ill{}} [anneau local noeth. complet,
\uncertain{sur} $S$] \struck{\ill{}}, et tout hom local $A \to B$,
$B$ fini sur $A$, et tt $\xi \in F(B)$, \struck{qui soit pro-modulaire}
[\ill{} que $\xi$ \uncertain{soit} modulaire], [on verra que ça implique :
pro-modulaire …] et tout $\mathfrak{p} \in \mathrm{Spec}(B)$,
$\xi_{\mathfrak{p}} \in F(B_{\mathfrak{p}})$, \uncertain{est} modulaire.
[Cf plus bas des conditions plus particulières)
\note{le numéro « 7° » est entouré ; « sans modulaires » semble écrit
au-dessus de « modulaire », lecture douteuse}

\marginal{[en fait, cela suffit, et cela signifie que
$\xi : \mathrm{Spec}\, B \to F$ est étale]}
\note{note dans la marge de gauche, à la hauteur de la fin de 7°}

\marginal{N.B. On \uncertain{pourrait} \ill{} que les conditions 1° à 5°,
[\uncertain{y compris} 6°] $+$ 7° impliquent l'analogue de 7° sans hyp. de
\uncertain{dim 1} \ill{}}
\marginal{Donner l'exemple type [$S$ de dim 1] montrant que (8°) n'est
pas conséquence des autres \ill{} $F$ \ill{} \ill{}. Est-ce que 6°
\uncertain{et} 7° \ill{} ? \emph{Non}, \ill{} \ill{} ce qui \ill{}
\uncertain{unique} \ill{} (\uncertain{par} \ill{}, \uncertain{continuité})}
\note{deux notes obliques dans la marge de gauche, séparées par un trait ;
la seconde est reliée à 7° et 8° par des flèches}

8°) [Pour tt \struck{préschéma} noeth. irréductible $S'$ sur $S$,
\struck{\ill{}} et $Z$ \uncertain{fini} sur $S'$,
$\xi \in F(Z) = F_{S'}(Z)$, $z \in Z$ au dessus du pt gén. de $S'$, tel
que $\xi$ soit modulaire
\note{le numéro « 8° » est entouré ; la phrase se poursuit en tête de la
page 15}

\page{14}
\emph{Immersion ouverte}

Conditions \struck{1° 2° 3° 4°} 1°, 2°, 3°, 4°, et de plus

a) $F \to S$ est un \emph{monomorphisme} \emph{étale} pour les arguments
artiniens.

\page{15}
en \uncertain{$Z$}, il existe une partie ouverte non vide $U$ de $Z$ telle que
$z' \in U$ implique que $\xi$ est \struck{étale en $z'$} modulaire en $z'$
[\uncertain{ou} simplement \emph{étale} en $z'$]
\note{le premier mot de la page achève la condition 8° de la page 13 ;
la lettre, lue $Z$, est celle de « $Z$ fini sur $S'$ », où l'on attend le
point $z$}

\emph{Corollaire 1} Pour que $F$ soit représentable par un $X$
\emph{étale} et \emph{séparé} sur $S$, il faut et il suffit qu'il
satisfasse les conditions \struck{1° à 6°} 1° 2° 3° 4° 6°, et de plus

\struck{(Et)} (Et) Le foncteur $F$ est étale (ou seulement, étale pour les
arguments artiniens) \struck{les $B$ de 5°) sont étales sur $A$}

\marginal{Expliciter \ill{} cas d'un \ill{} \ill{} ouvert}
\marginal{[dans ce cas plus direct et plus facile]}

\emph{Corollaire 2} Pour que $F$ soit représentable par un $X$ non
ramifié et séparé sur $S$, il faut et il suffit que les conditions 1° à
6°, \struck{8°} \struck{\ill{}} [ainsi que] soient vérifiées, et de plus
la condition 7°) [simplifiée :]
\note{« ainsi que » surcharge un mot biffé ; au-dessus, un « 8° » entouré
deux fois}

\marginal{7° \struck{\ill{}} (NR.)}

a) Le foncteur $F$ est non ramifié (i.e. les $B$ de 5°) sont non ramifiés
sur $A$ \struck{$A$})

b) Pour tout $A$ \struck{\ill{} \ill{}} de dim 1 \struck{\ill{}
\ill{}}, [irréductible] de $S$, \struck{\ill{} $A$ \ill{}} [\uncertain{idéal}
\ill{}] \struck{\ill{} $A$ \ill{}} \ill{} $\mathfrak{p}$, \struck{\ill{}}
$\forall$ \struck{\ill{}} [idéal] $I$ [tel que $I_{\mathfrak{p}} = 0$],
\struck{\ill{}} $\eta \in F(A/I)$, si
$\eta_{\mathfrak{p}} \in \mathrm{Im}\, F(A_{\mathfrak{p}})$, alors
$\eta \in \mathrm{Im}\, F(A)$.
\note{a) et b) sont réunis par une accolade ; la ligne de b) est deux fois
reprise, avec plusieurs passages biffés ; la lettre du point de
$F(A/I)$ est lue $\eta$}

\marginal{Expliquer 8° ici \uncertain{par} une impossibilité de
prolongements \ill{}}
\note{note oblique dans la marge de gauche, à la hauteur de b)}

\emph{Remarque 3} La condition 8°), est conséquence des autres
\ill{} \ill{} \uncertain{pour} que la conclusion \ill{} \ill{} :

\page{16}
(RG) Pour tt $S'$ \struck{\ill{}} [réduit-préschéma] fermé irréductible de
$S$, $\exists$ ouvert non vide $U \subset S'$ tel que $F_U$ soit
représentable par un préschéma loc. de type fini sur $U$.
\note{« RG » entouré, en tête de page}

\emph{Remarque 4} Limiter les anneaux $A$, $B$ et schémas $S'$ \ill{} qui
doivent être envisagés dans les conditions 2°, 4°, 5°, 7°, 8° ----

\emph{Exemples de foncteurs non ramifiés}

a) \ill{} \uncertain{triviaux}, plus généralement $G/H$, $H$ sous-groupe
\struck{ouvert} et \emph{fermé},

b) \struck{\ill{}} Classes de Correspondances divisorielles \ill{}
plat/$S$ …

c) Hom [\uncertain{foncteurs}] dans un groupe commutatif

b) Hom de V.A. dans un groupe comm. [Hom de V.A.]

e) Immersions [qui s'introduisent p.ex. dans les pbs de modules
\uncertain{primitifs} de \uncertain{diverses} \ill{}, p.ex. via problèmes
d'isogénies …] en particulier 1) morphismes diagonaux
$F \to F \times_S F$ 2) Foncteurs \uncertain{rectification}

d) Sous \struck{groupes} schémas, \ill{} d'un schéma en groupes
\note{une grande flèche courbe relie le « b) » de la quatrième ligne à
« d) » ; deux mots entourés, illisibles, sont écrits au-dessus de la fin
de d) et de f)}

f) Décomposition des fibres d'un morphisme en composantes connexes
(comp. en composantes irréd.). [Les foncteurs obtenus ici ne sont même
\emph{étales}]

Factorisation d'un morphisme fini \uncertain{mais ouvert} …

\emph{Applications en vue}. 1) Un schéma abélien sur $S$ loc. noeth.
\uncertain{connexe} est projectif sur $S$ [utiliser le schéma des
correspondances div. sur $X \times_S X$]

2) Existence du Picard d'un schéma propre sur un corps
\struck{[Seshadri-\ill{}]} [Chevalley-Seshadri-Grothendieck-Murre],
plus généralement, sur une base $S$, existence au-dessus d'un ouvert non
vide ----

\page{17}
Façon d'exprimer
\note{ces deux mots, en haut du feuillet, sont tout ce qu'il porte}

\page{19}
En effet, l'hyp. \ill{} qui est \uncertain{prouvée} dans la \ill{}
\ill{} \uncertain{remplie} : l'étude de $F \to F \times_S F$, \ill{}
\uncertain{prouve} que c'est \uncertain{repr.}, et grâce à 6°, c'est une
imm. fermée. Ceci posé, dans (A), $X \times_{F_{S'}} X$ est un sous
\struck{préschéma} \emph{fermé} [fini sur $X$,] de $X \times_S X$, qui
coïncide formellement avec la diagonale $X \to X \times_S X$, donc lui est
identique ; donc $X \to F$ est un \emph{monomorphisme}, d'autre part
étale, i.e. une imm. ouverte.

\emph{Dans} B, on prouve de façon analogue \ill{} $\exists\, U$ voisinage
ouvert de $x$, tel que $\xi | (X|U)$ \struck{est} soit un
\emph{monomorphisme}, comme \ill{} \ill{} \ill{} par « étale », \ill{}
\ill{}.
\note{cette page reprend les axiomes (A) et (B) de la page 8}

\end{document}
